\documentclass[11pt]{article}

\usepackage{amsmath,amssymb,amsthm,mathrsfs}
\usepackage{geometry}
\geometry{margin=1in}

\title{Kaprekar Dynamics as a Deterministic-to-Markov Lifted System\\
and Associated Inverse Problem Structure}
\author{A25 Framework (Computational Dynamics Note)}
\date{2026}

\newtheorem{definition}{Definition}
\newtheorem{theorem}{Theorem}
\newtheorem{lemma}{Lemma}

\begin{document}

\maketitle

\begin{abstract}
We study the Kaprekar 4-digit iteration as a deterministic dynamical system on a finite state space $\Omega^*$ and formalize its structure as a directed graph. We introduce a stochastic lifting procedure yielding a Markov chain representation, enabling spectral analysis via the normalized graph Laplacian. We further interpret the induced τ-depth distribution as a lossy observable of an ill-posed inverse problem over dynamical graphs. This framework separates deterministic collapse dynamics from stochastic relaxation geometry.
\end{abstract}

\section{Deterministic System}

\begin{definition}[State space]
Let
\[
\Omega^* = \{0,1,\dots,9999\} \setminus \{\text{repdigits}\},
\quad |\Omega^*| = 9990.
\]
\end{definition}

\begin{definition}[Kaprekar map]
Define
\[
K : \Omega^* \to \Omega^*
\]
by
\[
K(n) = \mathrm{desc}(n) - \mathrm{asc}(n),
\]
where digits are sorted in descending and ascending order.
\end{definition}

This defines a deterministic dynamical system $(\Omega^*, K)$.

\section{Graph Representation}

\begin{definition}[Transition graph]
Define a directed graph $G_K = (\Omega^*, E)$ with edges
\[
i \to K(i).
\]
\end{definition}

The deterministic transition operator is:
\[
P_{\mathrm{det}}(i,j) =
\begin{cases}
1 & j = K(i) \\
0 & \text{otherwise}.
\end{cases}
\]

This operator is not reversible and is highly non-invertible.

\section{Markov Lift}

\begin{definition}[Stochastic lift]
Let $U$ be the uniform transition kernel on $\Omega^*$. Define
\[
P_\varepsilon = (1-\varepsilon)P_{\mathrm{det}} + \varepsilon U,
\quad 0 < \varepsilon \le 1.
\]
\end{definition}

\begin{lemma}
For any $\varepsilon > 0$, $P_\varepsilon$ defines an irreducible and aperiodic Markov chain on $\Omega^*$.
\end{lemma}

\begin{proof}
Since $U$ has full support, every state communicates with every other state in one step with probability $\varepsilon / |\Omega^*| > 0$.
\end{proof}

\section{Markov Operator and Laplacian}

Let $\mathcal{P}$ be the Markov operator:
\[
(\mathcal{P} f)(x) = \sum_{y \in \Omega^*} P_\varepsilon(x,y) f(y).
\]

Let $\pi$ be the stationary distribution of $P_\varepsilon$.

\begin{definition}[Normalized Laplacian]
\[
L = I - \Pi^{1/2} P_\varepsilon \Pi^{-1/2}.
\]
\end{definition}

The eigenvalues of $L$ satisfy:
\[
0 = \lambda_0 \le \lambda_1 \le \cdots \le \lambda_{n-1}.
\]

The spectral gap is:
\[
\mu_1 = \lambda_1.
\]

\section{Hitting Time Structure}

\begin{definition}[Kaprekar depth]
Let $A = \{6174\}$. Define
\[
\tau(x) = \min \{t \ge 0 : K^t(x) = 6174\}.
\]
\end{definition}

This defines a deterministic hitting-time field over $\Omega^*$.

Under stochastic lift:
\[
\tau(x) = \mathbb{E}_P[T_A \mid X_0 = x],
\]
becomes a first-passage time of the Markov chain.

\section{Inverse Problem Formulation}

\begin{definition}[Inverse Kaprekar problem]
Given the observable distribution
\[
\{(\tau, N_\tau)\},
\]
reconstruct structural properties of the underlying graph $G_K$.
\end{definition}

This inverse map is ill-posed due to many-to-one collapse of $K$.

\begin{lemma}
The inverse mapping $\tau^{-1}$ is non-unique.
\end{lemma}

\begin{proof}
Multiple states share identical iteration depth under $K$, hence $\tau$ is not injective.
\end{proof}

\section{Regularized Reconstruction}

We define a variational reconstruction problem:
\[
\min_G \| \tau_G - \tau_{\mathrm{obs}} \|^2 + \lambda \mathcal{R}(G),
\]
where $\mathcal{R}(G)$ penalizes structural complexity.

This defines a discrete ill-posed inverse problem over dynamical graphs.

\section{Spectral Interpretation}

The spectral gap $\mu_1$ is not intrinsic to $K$ but emerges from the lifted operator $P_\varepsilon$.

\begin{theorem}[Spectral meaning]
$\mu_1$ governs the relaxation rate of probability mass toward the attractor under the Markov lift.
\end{theorem}

\begin{proof}
Standard Markov chain theory: $\mu_1$ controls convergence rate in $L^2(\pi)$.
\end{proof}

\section{Conclusion}

We separate three layers of structure:

\begin{itemize}
\item Deterministic collapse dynamics $(\Omega^*, K)$
\item Stochastic embedding $(\Omega^*, P_\varepsilon)$
\item Inverse problem geometry induced by $\tau$
\end{itemize}

Spectral properties arise only after stochastic lifting and should not be interpreted as intrinsic invariants of the deterministic map.

\end{document}
